% status: ZERO
% Chapter 11 of the second edition. Plan: docs/STRUCTURE.md. Rules: AUTHORING.md.
% Sources: research/FRONTIER.md (Part II additions, 2026-09-24).
\chapter{Below 16 Bits in Floating Point: FP8, MXFP4 and NVFP4}\label{ch:low-precision-float}

\begin{ChapterIntent}
Blackwell-class tensor cores multiply four-bit floating-point numbers
directly, provided each small block of values carries a scale. This chapter
explains the formats that make this work --- FP8's E4M3 and E5M2, the OCP
microscaling family with a shared power-of-two exponent per 32 values, and
NVFP4 with an FP8 scale per 16 values and a per-tensor scale --- and what each
trades between range, precision and overhead. It separates the tolerances of
weights, activations and the KV cache, and follows the shift from quantizing
trained models to training them in low precision, as in DeepSeek-V3's FP8
blocks and the MXFP4 checkpoints of gpt-oss and Kimi~K3.
\end{ChapterIntent}

\OpeningSection{117 billion parameters in 65 gigabytes}

OpenAI's gpt-oss-120b has 117~billion parameters \cite{openai2025gptoss}. In
the 16-bit format models are usually trained in, that would be 234~GB of
weights --- three 80~GB GPUs before a single token of cache. The checkpoint
OpenAI published is fourteen files totalling \textbf{65,248,893,184 bytes}
(measured: the sizes in the model repository's file listing, read 2026-09-24):
\textbf{4.46 bits per parameter}, and it runs on one 80~GB GPU.

Almost all of that saving comes from one decision: the mixture-of-experts
weights, about 115~billion of the 117~billion parameters, are stored in
MXFP4, a four-bit floating-point format in which every block of 32 numbers
shares one 8-bit scale. The attention weights, embeddings and router ---
about two billion parameters --- stay in 16 bits. The arithmetic checks
against the file: 114.7~billion expert parameters at 4.25 bits plus 2.1~billion
at 16 bits is 65.1~GB (derived from the published configuration; within 0.2\%
of the measured total).

Four-bit weights are not new; \Chref{weight-quantization} covers integer
formats that llama.cpp has used since 2023. What is new is that the four-bit
numbers are floating point, that the newest tensor cores multiply them
without converting them first, and that models are now \emph{trained} to be
served this way. This chapter explains the formats and what each one costs.

\NapkinSection{Napkin math: bits per value, including the scales}

A block-scaled format stores $k$ elements of $b_e$ bits and one scale of $b_s$
bits, so each value costs
\begin{equation}
b_{\text{eff}} = b_e + \frac{b_s}{k}\quad[\text{bits}].
\label{eq:bits-per-value}
\end{equation}
The formats this chapter covers, from their specifications
\cite{rouhani2023mx,nvidia2025nvfp4,micikevicius2022fp8}:

\begin{center}
\begin{tabular}{@{}lrrrr@{}}
\toprule
Format & Element & Block $k$ & Scale & Bits per value \\
\midrule
BF16 & E8M7 & --- & --- & 16 \\
FP8, per-tensor scale & E4M3 or E5M2 & tensor & FP32 & $\approx$8 \\
MXFP8 & E4M3 or E5M2 & 32 & E8M0 & 8.25 \\
MXFP6 & E2M3 or E3M2 & 32 & E8M0 & 6.25 \\
MXFP4 & E2M1 & 32 & E8M0 & 4.25 \\
NVFP4 & E2M1 & 16 & E4M3 (+ FP32 per tensor) & 4.5 \\
\bottomrule
\end{tabular}
\end{center}

Bytes per token follow directly: a dense 70-billion-parameter model reads about
140~GB per decode step in BF16, 70~GB in FP8, 37~GB in MXFP4 and 39~GB in NVFP4.
Compute follows the hardware. The B200's dense tensor-core peak doubles with
each halving of precision: 2.25~PFLOP/s in BF16, 4.5 in FP8, 9 in FP4 (NVIDIA
HGX specification, logged in \texttt{research/FRONTIER.md}). A format that
halves the bytes \emph{and} doubles the arithmetic keeps a workload's position
relative to the ridge (\Chref{roofline}) while making everything twice as
fast --- if the accuracy survives.

\section{Floating point with fewer bits}

A floating-point number spends its bits on a sign, an exponent that sets the
scale, and a mantissa that sets the precision within that scale. Take bits
away and one of the two suffers. The four-bit format used by both MXFP4 and
NVFP4, E2M1, has one sign bit, two exponent bits and one mantissa bit; its
non-negative values are 0, 0.5, 1, 1.5, 2, 3, 4 and 6
\cite{nvidia2025nvfp4}. They are not evenly spaced: the gap doubles as the
magnitude does, which matches the roughly bell-shaped distribution of weights
and activations better than an integer grid of the same size.

\EngineFigure{ch11-fp4-grid}{The eight non-negative values of FP4 (E2M1)
against a uniform 3-bit grid with the same maximum. Floating point spends its
resolution near zero, where most values lie; the integer grid spends it
evenly. A shared per-block scale decides which real interval these eight
points cover.}{fig:ch11-fp4-grid}

Eight magnitudes cannot represent a tensor on their own. The idea that makes
four-bit formats usable is the shared scale: a whole block of values is
divided by one scale factor chosen so that the block's largest magnitude lands
near the top of the grid. The element bits then only have to represent the
values' shape within the block; the scale carries the magnitude.

\section{FP8: E4M3 and E5M2}

The two FP8 formats standardized in 2022 split eight bits differently
\cite{micikevicius2022fp8}. E4M3 has four exponent bits and three mantissa
bits; it gives up infinities and keeps a single NaN pattern to extend its
range to $\pm448$. E5M2 has five exponent bits and two mantissa bits, keeps
IEEE-style special values, and reaches $\pm57{,}344$ with less precision. The
usual division of labour is E4M3 for weights and activations, where precision
matters more, and E5M2 for gradients in training, where range does --- but
DeepSeek-V3 used E4M3 everywhere, relying on fine-grained scales to supply the
range \cite{deepseek2024v3}.

\section{Scaling granularity: tensor, channel, block}

One scale per tensor is cheap and fragile: a single outlier anywhere in the
tensor sets the scale, and every other value is squeezed into a few levels of
the grid. One scale per row or column isolates outliers to their own channel.
One scale per small block isolates them to a few dozen values. Finer scaling
costs scale storage (equation~\eqref{eq:bits-per-value}) and, in a kernel,
the work of applying the scales.

DeepSeek-V3's FP8 recipe shows the choices a production system makes:
activations scaled per token per 128 channels ($1\times128$ tiles), weights per
$128\times128$ block \cite{deepseek2024v3}. Activations get the finer
granularity because they are computed at run time, carry outlier channels, and
are quantized on the fly; weights are quantized once, offline, and change
slowly across a block.

\section{Microscaling: shared exponents}

The OCP microscaling (MX) formats standardize block scaling for hardware: a
block of 32 elements shares one 8-bit scale in E8M0 --- an exponent with no
mantissa, so every scale is a power of two \cite{rouhani2023mx}. Multiplying by
a power of two is exact and cheap in hardware: it adjusts exponents without
rounding. The family covers MXFP8, MXFP6, MXFP4 and MXINT8, each with the same
32-element blocks. Because the scale can only be a power of two, a block's
largest value may land anywhere between the top half and the top of the
element grid, wasting up to one bit of its resolution.

\section{NVFP4: two-level scaling}

NVIDIA's NVFP4 keeps the E2M1 elements but changes the scaling
\cite{nvidia2025nvfp4}: blocks of 16 instead of 32, and a scale stored in E4M3
--- a real fraction, not just a power of two --- chosen to fit the block
closely. Because an E4M3 scale has limited range of its own, a second, FP32
scale per tensor brings each tensor's scales into the E4M3 range. The price
is a quarter of a bit per value more than MXFP4 (4.5 against 4.25); the gain is
smaller blocks and exact-fit scales. NVIDIA reports that DeepSeek-R1-0528
served in NVFP4 lost 1\% or less of accuracy against FP8 on its evaluations
(vendor measurement, not reproduced here).

\section{Tensor cores that consume scales}

A format is only fast if the hardware multiplies it directly. Blackwell's
fifth-generation tensor-core instructions include block-scaled variants that
take the scale factors alongside the operands: one UE8M0 scale per 32 elements
of the reduction dimension for the MX types, one UE4M3 scale per 16 for NVFP4
(CUTLASS documentation, logged in \texttt{research/FRONTIER.md}). The kernel
never expands a block to 16 bits in registers; the tensor core applies the
scales inside the multiply--accumulate.

Hopper has FP8 tensor cores but no FP4, and its FP8 path has a subtlety that
DeepSeek found in production training: on H800 GPUs, FP8 matrix products keep
only about 14 bits of precision in their accumulators. Summing thousands of
products at that precision loses accuracy, so DeepSeek copies partial sums to
FP32 registers every 128 elements of the reduction and accumulates there
\cite{deepseek2024v3}. The same arithmetic happens in an FP8 inference GEMM on
the same hardware; the fix costs some of the speed the format bought.

\section{Weights, activations and the cache}

The three kinds of tensor tolerate low precision differently. Weights are
fixed: they can be quantized once, with care, offline, and checked. Activations
are computed per request, contain outlier channels, and must be quantized on
the fly with scales computed from the data itself, which costs a reduction per
block at run time. The KV cache is written once and read at every later step;
its errors enter every attention score that follows
(\Chref{kv-quantization}). A deployment chooses per tensor: gpt-oss keeps
attention and embeddings in 16 bits and stores only expert weights in MXFP4;
Kimi~K3 reports MXFP4 weights with MXFP8 activations
\cite{moonshot2026kimik3}.

\section{Training in low precision changes inference}

Until recently, low precision was something done to a model after training,
with an accuracy cost to be measured and minimized. The 2025--26 checkpoints
changed that. DeepSeek-V3 was trained in FP8 and published with FP8 weights in
$128\times128$ blocks, so serving it in FP8 introduces no conversion at all
\cite{deepseek2024v3}. gpt-oss was post-trained with its expert weights in
MXFP4, and its reported evaluations use the quantized model
\cite{openai2025gptoss}. Kimi~K3 reports quantization-aware training for its
MXFP4 weights \cite{moonshot2026kimik3}. NVIDIA trained a 12-billion-parameter
model on 10~trillion tokens in NVFP4 with loss matching an FP8 baseline
\cite{nvidia2025nvfp4pretrain}. When a model is trained or post-trained in a
format, the question ``how much accuracy does quantization cost?'' changes to
``what does it cost to serve this model in any \emph{other} format?''

\LedgerSection

Serving a model's large matrices in a four-bit block-scaled format (MXFP4 or
NVFP4) instead of BF16:

\begin{ConstraintLedger}
Memory bandwidth & buys & 3.6--3.8$\times$ fewer bytes per weight read (equation~\eqref{eq:bits-per-value}). \\
Memory capacity & buys & gpt-oss-120b in 65~GB instead of about 234~GB (measured file sizes; derived). \\
Compute & buys & 4$\times$ the dense tensor-core peak of BF16 on Blackwell; nothing on Hopper, which has no FP4 tensor cores. \\
Correctness & risks & Block-scale error and outliers; negligible when the model was trained for the format, to be measured when it was not. \\
\end{ConstraintLedger}

\LabSection{round-trip three formats and compare their errors}

\LabIntent{In \texttt{code/labs/ch11-block-formats/}, reference encoders and
decoders for E4M3, MXFP4 (32-element blocks, power-of-two E8M0 scales) and
NVFP4 (16-element blocks, E4M3 scales, FP32 tensor scale), written from the
specifications and tested against every representable value. The lab
round-trips real weight rows --- dequantized from a GGUF file --- and synthetic
rows with injected outliers, and reports relative error per format and block
size. Break it: put one large outlier in a block and watch which formats
lose the block's small values.}

\HappenedSection{fourteen bits of accumulator}

DeepSeek's FP8 training run is the clearest recorded case of a low-precision
format failing silently in a production kernel and being fixed by design
\cite{deepseek2024v3}. FP8 inputs were not the problem; the accumulator was.
The H800's FP8 matrix instructions accumulate at about 14 bits of precision,
enough for short reductions and not for the thousands of products in a
large matrix product. The fix changed the kernel's structure rather than the
format: partial sums leave the tensor cores every 128 elements for FP32
registers on the ordinary cores, where the long sum is completed. It is a
reminder that a format's specification describes its storage, not the
precision of every operation performed on it --- and that the only way to find
the difference is to test against a higher-precision oracle.

\FrontierSection{2026-09-24}

Four-bit floating point moved from research to shipping checkpoints in about a
year: MXFP4 in gpt-oss (August 2025) and Kimi~K3 (2026), NVFP4 pretraining
at 10~trillion tokens \cite{nvidia2025nvfp4pretrain}. Blackwell Ultra parts
list more dense FP4 throughput than B200 while trimming INT8 and FP64 (NVIDIA
HGX specification), a sign of where hardware designers expect inference
arithmetic to go. AMD's MI355X goes further for six bits: it lists 10.1
dense PFLOP/s for MXFP6 and MXFP4 alike, twice its FP8 rate and four times its
BF16 rate, so on that part a six-bit format costs no more arithmetic time than
a four-bit one \cite{amd2025mi355x}. \todo{verify: Rubin-class FP4
specifications from NVIDIA}.

\ExercisesSection

\begin{enumerate}
\item \emph{Derivation.} Using equation~\eqref{eq:bits-per-value}, compare
MXFP4 and NVFP4 for a 1-trillion-parameter model. How many H200s does each
need for weights alone?
\item \emph{Prediction.} A block of 32 values has one value of 100 and the
rest between $-1$ and 1. Predict which of MXFP4 and NVFP4 represents the small
values better, and by how much.
\item \emph{Implementation.} Extend the lab to MXFP6 and measure where it sits
between MXFP4 and MXFP8 in error per bit.
\item \emph{Falsification.} Find a matrix for which per-tensor FP8 beats
per-block MXFP8 in error. What property must it have?
\end{enumerate}
